

Šajā promocijas darbā pētīta \Gls{snn} izmantošana efektīvai telpisku un temporālu datu apstrādei, galvenokārt pievēršoties to paātrināšanai pārkonfigurējamā aparatūrā ar \Gls{fpga}. Darbs aptver gan algoritmiskos, gan aparatūras aspektus \Gls{snn} realizācijā, piemēram, neironu modeļus, neironu tīklus, apmācības metodes, telpiski temporālu datu reprezentācijas un \Gls{fpga} paātrinātāju arhitektūras. Rezultāti demonstrē, ka \Gls{snn} nodrošina piemērotu skaitļošanas pieeju retinātiem temporāliem datiem un tos var realizēt pārkonfigurējamā aparatūrā, lai sasniegtu zema latentuma un energoefektīvu inferenci.

Pirmajā tēzē pārbaudīta \Gls{fpga} izmantošana kā pārkonfigurējama aparatūras platforma \Gls{snn} paātrināšanai. Izstrādāts pārkonfigurējams, uz PCIe balstīts \Gls{fpga} paātrinātājs konvolucionālo \Gls{snn} realizācijai. Paātrinātājs atbalsta konfigurējamus konvolūcijas slāņus un impulsu neironu modeļus dažādu \Gls{snn} arhitektūru realizācijai aparatūras platformā bez atkārtotas sintēzes. Pārbaudīts, ka kvantēšana ar mazu bitu platumu var saglabāt konvolucionālo \Gls{snn} klasifikācijas veiktspēju, vienlaikus ļaujot tos realizēt pārkonfigurējamā \Gls{fpga} paātrinātājā. Rezultāti parādīja, ka svaru precizitātes samazināšana no 32 līdz 4 bitiem saglabāja F1 rādītāju 1~\% robežās no ekvivalenta pilnas precizitātes \Gls{snn} modeļa gan \textit{MNIST}, gan četrslāņu un sešslāņu \textit{Caltech-101} modeļiem. Konkrēti, attiecīgie 4 bitu modeļi sasniedza F1 rādītājus 96,79~\%, 89,32~\% un 88,10~\%. Novērtētā realizācija demonstrēja arī mazu dinamisko jaudas patēriņu, apliecinot energoefektīvas \Gls{snn} klasifikācijas iespējamību pārkonfigurējamā aparatūrā. Eksperimentālie rezultāti ar \textit{Kintex UltraScale KU115} \Gls{fpga} parādīja, ka izmērītais vidējās dinamiskās jaudas pieaugums \Gls{fpga} inferences laikā pārbaudītajā slodzē sasniedza aptuveni 91~mW. Jaudas mērījumi arī uzsver atšķirību starp paātrinātāja aktivitāti un aparatūras platformas dīkstāves jaudas patēriņu. Vidējā jauda, kas izmērīta 12~V barošanas ķēdē dīkstāvē, bija 15,77~W. Salīdzinoši nelielais inferences jaudas pieaugums norāda, ka \Gls{snn} slodze radīja ierobežotu papildu vidējo jaudu. Salīdzinoši lielā dīkstāves jauda izskaidrojama ar lieljaudas KU115 \Gls{fpga} izstrādes kartes izmantošanu jaudas ziņā optimizētas iegultās platformas vietā. Specializēta paātrinātāja platforma ar piemērota izmēra \Gls{fpga} samazinātu šo dīkstāves patēriņu. Tomēr šādi uzlabojumi paliek turpmākā darba virziens.

Eksperimenti demonstrēja arī zema latentuma inferenci novērtētajiem konvolucionālajiem tīkliem un parādīja, ka \Gls{fpga} paātrināšana var nodrošināt determinētu konveijerizētu izpildi un latentuma uzlabojumus salīdzinājumā ar iegultām GPU platformām (piemēram, \textit{NVIDIA Jetson AGX Orin}). Eksperimenti arī demonstrēja, ka konveijerizētas \Gls{fpga} izpildes priekšrocības ir izteiktākas, ja kaskādē savienoti vairāki konvolūcijas slāņi. Salīdzinājums ar augstas veiktspējas diskrētu GPU (piemēram, \textit{NVIDIA GeForce RTX 5070 Ti}) parādīja, ka \Gls{fpga} paātrināšana ne vienmēr nodrošina visaugstāko caurlaidību. Šajā gadījumā galvenās priekšrocības ir determinēts latentums un ievērojami mazākas jaudas prasības nekā GPU platformām.

Otrajā tēzē pārbaudīta konvolucionālo \Gls{snn} piemērotība objektu detektēšanai punktu mākoņos un novērtēta to potenciālā energoefektivitāte. Eksperimentālie rezultāti demonstrēja, ka konvolucionālie \Gls{snn} spēj apstrādāt laikā kodētas 3D telpiskās reprezentācijas. Konkrēti, izstrādātā 3D-SCNN arhitektūra piemērota objektu detektēšanai LiDAR datos. Izstrādātā impulsu detektēšanas arhitektūra apvieno temporālo vokseļu kodēšanu, impulsu \textit{VGG} pamattīklu un detektēšanas galvas objektu klasifikācijai un orientētu ietverošo taisnstūru regresijai. \textit{KITTI} validācijas datu kopā \textit{VGG-13} modelis sasniedza 72,03~\% \Gls{bev} vidējo precizitāti automobiļu klasei. Papildu eksperimenti ar izkliedētu \Gls{tof} sensoru sistēmu demonstrēja praktisku piemērotību \Gls{tof} sensoru punktu mākoņu datiem. Analītisks enerģijas salīdzinājums, kas balstīts uz reprezentatīviem 45~nm CMOS operāciju enerģijas pieņēmumiem, novērtētajām \Gls{snn} arhitektūrām uzrādīja aptuveni par 27,5~\% mazāku inferences enerģiju nekā ekvivalentiem tradicionāliem \Gls{ann} modeļiem. Eksperimenti demonstrēja, ka uz \Gls{snn} balstīti modeļi var iegūt telpiskas un temporālas pazīmes no retinātiem 3D sensoru datiem, kas kodēti uz notikumiem balstītā reprezentācijā, lai veiktu objektu detektēšanu un izmantotu energoefektīvas notikumu apstrādes priekšrocības.

Tomēr piedāvātajam 3D-SCNN ir noteikti ierobežojumi. Parādīts, ka \Gls{snn} efektivitātes priekšrocības ir stipri atkarīgas no impulsu retinājuma un simulācijas laika soļu skaita. Atšķirībā no tradicionālajiem \Gls{ann}, kuros aritmētisko operāciju skaitu nosaka tīkla arhitektūra, \Gls{snn} operāciju skaits papildus ir atkarīgs no ģenerēto impulsu skaita. Tādējādi reta impulsu aktivitāte var būtiski samazināt operāciju skaitu. Tomēr rezultāti arī parāda, ka retinātu aprēķinu ieguvumu var kompensēt atmiņas piekļuves, īpaši tad, ja neironu stāvokļi tiek glabāti ārējā SDRAM atmiņā. Tādēļ efektīva atmiņas shēma un mikroshēmas iekšējās \Gls{bram} izmantošana ir būtiska \Gls{snn} potenciālo priekšrocību īstenošanai.

Turklāt piedāvātās 3D-SCNN detektēšanas arhitektūras novērtētie enerģijas ietaupījumi vēl nav pārbaudīti ar pilnībā realizēta modeļa fiziskiem jaudas mērījumiem. Turpmākajā darbā joprojām nepieciešama impulsu pamattīkla, pazīmju piramīdas starpposma un detektēšanas galvas posmu pilnīga realizācija aparatūras platformā, kam sekotu jaudas patēriņa mērījumi.

Trešajā tēzē pārbaudīta konveijerizēta, pārkonfigurējama \Gls{fpga} paātrinātāja iespējamība ar jauktu atmiņas shēmu, kur svari tiek glabāti mikroshēmā un neironu membrānas potenciāli — ārpus tās. Šī izvēle risina ierobežotas iekšējās atmiņas problēmu, realizējot dziļākus konvolucionālos \Gls{snn}, vienlaikus saglabājot elastību dažādām tīklu konfigurācijām. Novērtētajiem četrslāņu un sešslāņu \textit{Caltech-101} tīkliem novērtētās mikroshēmas iekšējās atmiņas prasības samazinātas attiecīgi aptuveni par 53,8~\% un 55,0~\% salīdzinājumā ar ekvivalentām konfigurācijām, kur visa atmiņa izvietota mikroshēmā. Turklāt noteikts, ka novērtētie četrslāņu un sešslāņu tīkli sasniedza attiecīgi 2,43~ms un 2,81~ms inferences latentumu, parādot, ka membrānas potenciālu atmiņas pārvietošana ārpus mikroshēmas saglabā zemu inferences latentumu. Papildu eksperimenti demonstrēja arī konveijerizētas izpildes priekšrocības dziļākām tīkla konfigurācijām un apstiprināja arhitektūras piemērotību zema latentuma un determinētai \Gls{snn} inferencei.

Kopumā rezultāti apstiprina promocijas darba mērķa sasniegšanu un demonstrē, ka kvantētu \Gls{snn} modeļu inferenci un membrānas potenciālu ārēju glabāšanu var apvienot efektīvai resursu izmantošanai. Turklāt piedāvātais risinājums līdzsvaro pārkonfigurējamību, mērogojamību un energoefektivitāti uz \Gls{fpga} balstītā neiromorfajā skaitļošanā. Uz notikumiem balstītu \Gls{snn} realizācija retinātu telpiski temporālu datu apstrādei un to ieviešana pārkonfigurējamā aparatūrā nodrošina īstenojamu pieeju mazas jaudas inferencei perifērijas ierīcēs. Tomēr saglabājas vairāki ierobežojumi. Jo īpaši novērtēto punktu mākoņu \Gls{snn} modeļu detektēšanas precizitāte ir zemāka nekā tradicionālajiem \Gls{ann} modeļiem ar sarežģītākām arhitektūrām. Arī \Gls{snn} efektivitātes priekšrocības ir atkarīgas no vairākiem faktoriem, piemēram, impulsu retinājuma, kodēšanas laika soļa ilguma, kvantēšanas un atmiņas arhitektūras. Tādēļ efektīvai \Gls{snn} inferencei nepieciešama kopīga neironu tīkla modeļa, datu reprezentācijas un aparatūras paātrinātāja realizācijas optimizācija.
